%% kapa: exam=1 round=22 year=2011 period=1 form=A subject=civil_law subject_ko=민법 seq=3 number=3
%% source: Q-Net 감정평가사 기출문제 (cst00309) artlSeq=5036790 file=extracted/1008/5036790/2011년 제22회 감정평가사 1차시험 기출문제/1교시.hwp
%% conversion: kordoc 4.24.51 -> tools/gen1.py (auto); review=proofread-sonnet; flags=none
\begin{question}{3}
민법상 법인에 관한 설명으로 옳지 않은 것은? (다툼이 있으면 판례에 의함)
\begin{choices}
  \item 법인은 대표권이 없는 이사의 행위에 대해서 ‘법인의 불법행위책임’(제35조)을 진다.
  \item 법인과 이사의 이익이 상반하는 사항에 관해서 그 이사는 대표권이 없다.
  \item 이사가 대표권 없이 한 법률행위에는 민법의 무권대리 또는 표현대리의 규정이 준용된다.
  \item 이사의 대표권은 정관으로 제한할 수 있고, 이를 등기하면 제3자에게 대항할 수 있다.
  \item 이사는 그 임기가 만료되거나 이사직을 사임한 경우에도 후임이사가 선임될 때까지는 계속해서 이사의 직무를 수행할 수 있는 업무수행권을 갖는다.
\end{choices}
\end{question}
\answer{1}
